% Claim proof file for row claim_3_12 -- "HardInterventionNodeSplitOrder".
% Auto-generated stub by scaffold/solving_current_row.py at row start. A
% manager agent dedicated to writing the TeX proof fills in the body below.
% The proof must:
%   - Stay close to the lecture notes' proof if one exists (always search
%     the LN first for a `\begin{proof}` block immediately following the
%     claim; copy / lightly edit when present).
%   - Otherwise use the LN paradigm and cite prior claims by their `ref`.
%   - Be self-contained enough that a mathematician can follow it without
%     re-reading the LN.
%   - Have no `sorry`, `omitted`, or "obvious" shortcuts.
%
% Compiles standalone via the `subfiles` package.

\documentclass[main]{subfiles}

\begin{document}
% ref: claim_3_12 (proof, with statement restated for readability)
% title: HardInterventionNodeSplitOrder
\def\rowref{claim\_3\_12}
\def\rowtitle{HardInterventionNodeSplitOrder}
\phantomsection\label{claim_3_12}
%
% Cross-references: see claim_<ref>_statement_<title>.tex in this folder
% for the corresponding statement.

% --- statement (restated) ---------------------------------------------------
\begin{Rem}
    Note that if $W_1$ and $W_2$ are not disjoint and $w \in W_1 \cap W_2 \ins V$
    then first hard intervening on $w$ turns $w$ into an input node, for now indicated as $w^i$, and a node-splitting hard intervention (if we would define it for input nodes) would not change $w^i$.
    If, on the other hand, we would first split the node $w$ into $w^o$ and $w^i$ then we would first need to resolve the ambiguity on which of those two the hard intervention should be applied. A hard intervention on $w^i$ would not do anything, but would leave the additional output node $w^o$ in the graph, while hard intervening on $w^o$ would turn $w^o$ into an additional input node, for now indicated as $(w^o)^i$. So in the latter case we are left with two input node $(w^o)^i$, which does not have any edges, and $w^i$, which might have outgoing edges.
\end{Rem}

% --- proof ------------------------------------------------------------------
\begin{proof}
% TODO: write the proof body.
\end{proof}

\end{document}
